\documentclass[tikz,border=8pt]{standalone}
\usepackage{ctex,amsmath,amssymb}
\usetikzlibrary{arrows.meta,positioning,calc,fit,backgrounds}
% Shared visual language for LFS architecture diagrams.
\RequirePackage{../../mymath}

\definecolor{LFSBlue}{HTML}{2563EB}
\definecolor{LFSBlueSoft}{HTML}{DBEAFE}
\definecolor{LFSGreen}{HTML}{15803D}
\definecolor{LFSGreenSoft}{HTML}{DCFCE7}
\definecolor{LFSOrange}{HTML}{C2410C}
\definecolor{LFSOrangeSoft}{HTML}{FFEDD5}
\definecolor{LFSRed}{HTML}{B91C1C}
\definecolor{LFSRedSoft}{HTML}{FEE2E2}
\definecolor{LFSPurple}{HTML}{7E22CE}
\definecolor{LFSPurpleSoft}{HTML}{F3E8FF}
\definecolor{LFSGray}{HTML}{475569}
\definecolor{LFSGraySoft}{HTML}{F1F5F9}

\tikzset{
  lfs picture/.style={
    font=\small,
    node distance=8mm and 12mm,
    >=Latex,
    every path/.style={line cap=round, line join=round}
  },
  block/.style={
    draw=LFSBlue,
    fill=LFSBlueSoft,
    rounded corners=2pt,
    line width=0.8pt,
    align=center,
    inner xsep=7pt,
    inner ysep=5pt,
    minimum height=10mm
  },
  compute/.style={
    block,
    draw=LFSPurple,
    fill=LFSPurpleSoft
  },
  storage/.style={
    block,
    draw=LFSGreen,
    fill=LFSGreenSoft
  },
  interface/.style={
    block,
    draw=LFSOrange,
    fill=LFSOrangeSoft
  },
  risk/.style={
    block,
    draw=LFSRed,
    fill=LFSRedSoft
  },
  neutral/.style={
    block,
    draw=LFSGray,
    fill=LFSGraySoft
  },
  decision/.style={
    diamond,
    draw=LFSOrange,
    fill=LFSOrangeSoft,
    line width=0.8pt,
    align=center,
    inner sep=2pt,
    aspect=2
  },
  group/.style={
    draw=LFSGray,
    rounded corners=3pt,
    dashed,
    line width=0.7pt,
    inner sep=6pt
  },
  arrow/.style={
    -{Latex[length=2.4mm,width=1.6mm]},
    draw=LFSGray,
    line width=0.9pt
  },
  data arrow/.style={
    arrow,
    draw=LFSBlue
  },
  control arrow/.style={
    arrow,
    draw=LFSOrange,
    dashed
  },
  residual/.style={
    arrow,
    draw=LFSGreen,
    rounded corners=4pt
  },
  label text/.style={
    font=\scriptsize,
    text=LFSGray,
    fill=white,
    inner sep=1.5pt
  },
  title/.style={
    font=\bfseries,
    text=LFSGray,
    align=center
  }
}


\begin{document}
\begin{tikzpicture}[my picture, node distance=5mm and 6mm]

% 5层技术栈卡片定义
% 第五层：应用智能体
\node[interface, text width=128mm, minimum height=11mm] (l5) {
  \textbf{应用智能体系统（第二篇）}\\[1mm]
  \scriptsize 检索增强生成（RAG） \quad$\bullet$\quad 工具调用与沙箱执行 \quad$\bullet$\quad 任务规划与记忆状态 \quad$\bullet$\quad 垂直领域落地闭环
};

% 第四层：服务调度
\node[compute, text width=128mm, minimum height=11mm, below=of l5] (l4) {
  \textbf{模型服务与推理调度（第四篇）}\\[1mm]
  \scriptsize 连续批处理（Continuous Batching） \quad$\bullet$\quad 分页增量 KV Cache \quad$\bullet$\quad Prefill/Decode 分离 \quad$\bullet$\quad 投机解码与算子融合
};

% 第三层：训练对齐
\node[storage, text width=128mm, minimum height=11mm, below=of l4] (l3) {
  \textbf{模型训练与对齐方法（第三篇）}\\[1mm]
  \scriptsize 预训练规模定律（Scaling Laws） \quad$\bullet$\quad 监督微调（SFT） \quad$\bullet$\quad 参数高效微调（PEFT/LoRA） \quad$\bullet$\quad 人类偏好对齐（PPO/DPO/GRPO）
};

% 第二层：模型架构
\node[block, text width=128mm, minimum height=11mm, below=of l3] (l2) {
  \textbf{模型架构与表示基础（第一篇 / 第五篇）}\\[1mm]
  \scriptsize 词元化与向量嵌入 \quad$\bullet$\quad 自注意力与 Transformer \quad$\bullet$\quad 混合专家模型（MoE） \quad$\bullet$\quad 旋转位置编码与长上下文
};

% 第一层：基础设施
\node[neutral, text width=128mm, minimum height=11mm, below=of l2] (l1) {
  \textbf{算力与数据基础设施（第四篇 / 第三篇）}\\[1mm]
  \scriptsize AI 加速芯片与高带宽显存（HBM） \quad$\bullet$\quad 高速集合通信网络 \quad$\bullet$\quad 3D 混合分布式并行 \quad$\bullet$\quad 海量语料清洗流水线
};

% 层间流动箭头
\draw[data arrow] ($(l1.north)+(0,0.5mm)$) -- ($(l2.south)+(0,-0.5mm)$);
\draw[data arrow] ($(l2.north)+(0,0.5mm)$) -- ($(l3.south)+(0,-0.5mm)$);
\draw[data arrow] ($(l3.north)+(0,0.5mm)$) -- ($(l4.south)+(0,-0.5mm)$);
\draw[data arrow] ($(l4.north)+(0,0.5mm)$) -- ($(l5.south)+(0,-0.5mm)$);

% 侧边标注
\node[label text, rotate=90, left=4mm of $(l1)!0.5!(l2)$] {底座支撑};
\node[label text, rotate=90, left=4mm of $(l3)!0.5!(l4)$] {工程落地};
\node[label text, rotate=90, left=4mm of l5] {业务赋能};

\end{tikzpicture}
\end{document}
